% ================================================================
% SECTION III: PROTOCOL ARCHITECTURE COMPARISON
% ================================================================

We compare the six protocols on encoding and redundancy, commitment scheme, and verification paradigm.
These choices fix a protocol's bandwidth cost, what a light client can check on its own, and which parties must be trusted.

% ----------------------------------------------------------
% III-A. ENCODING AND REDUNDANCY
% ----------------------------------------------------------
\subsection{Encoding and Redundancy}
\label{subsec:encoding}

All six protocols use Reed--Solomon (RS) erasure coding, but they differ in dimensionality, redundancy factor and how data is reconstructed.
Table~\ref{tab:encoding_comparison} summarizes these differences.

\begin{table}[t]
\centering
\caption{Encoding and redundancy.}
\label{tab:encoding_comparison}
\setlength{\tabcolsep}{2pt}
\footnotesize
\begin{tabular}{l l c l}
\toprule
\textbf{Protocol} & \textbf{Encoding} & \textbf{Redund.} & \textbf{Reconstruction} \\
\midrule
Polkadot~\cite{Jeff_Burdges_ELVES} & 1D RS (validators) & ${\sim}3\times$ & $f{+}1$ of $n$ ($n{=}3f{+}k$) \\
Ethereum~\cite{EIP-7594} & 1D RS (PeerDAS cells) & ${\sim}2\times$ & 50\%+ of 128 columns \\
Celestia~\cite{AlBassam_LazyLedger} & 2D RS ($k \times k \to 2k \times 2k$) & ${\sim}4\times$ & $k{\times}k$ shares \\
Espresso~\cite{Espresso:Hotshot_and_Triamisu} & Polynomial VID & $1/r$, $r{=}m/n$ & $m$ of $n$ evaluations \\
NEAR~\cite{NEAR_Nightshade_Scaling_Blockchain} & 1D RS (chunk parts) & ${\sim}3\times$ & $1/3$ of parts \\
Avail~\cite{Avail_The_DA_Blockchain} & 1D RS (row-wise) & ${\sim}2\times$ & $n$ of $2n$ coded pieces \\
\bottomrule
\end{tabular}
\end{table}

Higher redundancy costs proportionally more bandwidth and aggregate storage; in return, data can be recovered when a larger fraction of coded pieces is missing.

Celestia has the highest redundancy at ${\sim}4\times$ due to its two-dimensional RS structure: original data of dimension $k \times k$ is extended to $2k \times 2k$, yielding four times the original data volume in shares~\cite{AlBassam_LazyLedger}.
The overhead pays for two-dimensional sampling, in which light clients sample from both the rows and the columns of the extended matrix.
The 2D layout also allows a single row or column of length $2k$ to be recovered from any $k$ of its shares, so a node can repair one row without reconstructing the whole block.

Polkadot and NEAR both use ${\sim}3\times$ redundancy, distributing erasure-coded pieces across their respective validator sets.
In Polkadot's ELVES protocol, parachain block data (Proofs of Validity, or PoVs) are RS-encoded and distributed such that each validator stores one coded chunk, with reconstruction requiring $f{+}1$ of $n$ chunks where $n = 3f + k$~\cite{Jeff_Burdges_ELVES, Overview_of_polkadot}.
NEAR's Nightshade sharding similarly encodes chunk data with ${\sim}3\times$ redundancy, requiring only $1/3$ of the parts for reconstruction~\cite{NEAR_Nightshade_Scaling_Blockchain}.
In both, responsibility for the data is spread over the validator set; there is no separate committee and no sampling.

Ethereum and Avail have the lowest redundancy, ${\sim}2\times$.
Ethereum's EIP-4844 blob transactions encode data into 128-byte cells arranged across 128 columns. PeerDAS extends each blob's polynomial evaluation points such that any 64 of the 128 columns suffice for reconstruction~\cite{EIP-7594}.
Avail performs row-wise RS encoding on its application data matrix, extending $n$ original pieces to $2n$ coded pieces per row, so that any $n$ of the $2n$ pieces suffice for recovery~\cite{Avail_The_DA_Blockchain}.
With a $2\times$ extension any half of the coded pieces recovers the data, which together with KZG commitments is enough for DAS.

Espresso's Tiramisu uses a polynomial-based verifiable information dispersal (VID) scheme.
Data is represented as a polynomial that is evaluated at $n$ distinct points (one per committee member); any $m$ evaluations suffice for interpolation, where the code rate $r = m/n$ is configurable~\cite{Espresso:Hotshot_and_Triamisu}.
Algebraically this is an RS code, but because it is built into the VID commitment layer, one proof certifies both that the encoding is correct and that each member received its piece.

% ----------------------------------------------------------
% III-B. COMMITMENT SCHEMES AND PROOF SYSTEMS
% ----------------------------------------------------------
\subsection{Commitment Schemes and Proof Systems}
\label{subsec:commitments}

The commitment scheme determines what a light client can verify, whether a trusted setup is needed, and the cost of producing and checking proofs.
The six protocols use three families: KZG polynomial commitments, Namespaced Merkle Trees (NMTs) and standard Merkle trees.
Table~\ref{tab:commitment_comparison} compares their properties.

\begin{table}[t]
\centering
\caption{Commitment scheme properties (KZG: Avail, Ethereum, Espresso; NMT: Celestia; Merkle: Polkadot, NEAR).}
\label{tab:commitment_comparison}
\setlength{\tabcolsep}{2pt}
\footnotesize
\begin{tabular}{l c c c}
\toprule
\textbf{Property} & \textbf{KZG} & \textbf{NMT} & \textbf{Merkle} \\
\midrule
Proof size & $O(1)$, ${\sim}48$ bytes & $O(\log n)$ & $O(\log n)$ \\
Trusted setup & Yes (SRS) & No & No \\
Namespace filter & No & Native & No \\
Verification & Pairings & Hashes & Hashes \\
Encoding correctness & Algebraic & Fraud proofs & Not verifiable \\
\bottomrule
\end{tabular}
\end{table}

KZG polynomial commitments~\cite{KZG_Original} are used by Avail, Ethereum, and Espresso.
A KZG commitment to a polynomial $p(x)$ of degree $d$ is a single elliptic curve group element (${\sim}48$ bytes on BLS12-381), and an opening proof for any evaluation $p(z) = v$ is also a single group element.
Verification requires two pairing operations, which are computationally more expensive than hash evaluations but yield $O(1)$-size proofs regardless of the data size.
KZG commitments enable \emph{algebraic verification of encoding correctness}: given commitments to the rows (or columns) of an erasure-coded data matrix, anyone can verify that the extended portions are consistent with the original data by checking polynomial relationships between commitments, without accessing the underlying data.
This lets Avail and Ethereum light clients reject incorrectly encoded data without fraud proofs.
The drawback is that KZG needs a structured reference string (SRS) from a trusted setup; anyone holding the setup's secret could forge proofs.
Ethereum and Avail both ran multi-party ceremonies to reduce this risk, Ethereum's with over 140{,}000 contributors.

Celestia uses Namespaced Merkle Trees (NMTs)~\cite{AlBassam_LazyLedger}.
An NMT is a standard Merkle tree augmented with namespace identifiers: each leaf is tagged with a namespace, internal nodes carry the minimum and maximum namespace of their subtrees, and the tree is sorted by namespace.
This structure enables \emph{namespace inclusion and exclusion proofs}. A light client can request all shares belonging to a specific namespace (e.g., a particular rollup's data) and receive a Merkle proof that the returned set is complete (no shares from that namespace have been omitted, and no extraneous shares have been included).
Rollup nodes use this to download only their own data instead of the whole block.
NMT proofs are $O(\log n)$ in size.
Encoding correctness, however, is \emph{not} algebraically verifiable from the NMT alone. Instead, Celestia relies on Bad Encoding Fraud Proofs (BEFPs) to detect and propagate evidence of incorrect erasure coding~\cite{AlBassam_FraudDA}.

Polkadot and NEAR use standard Merkle trees.
Polkadot commits to erasure-coded PoV chunks via a Merkle trie whose root is included in the relay chain block header~\cite{Overview_of_polkadot, polkadot_article}.
NEAR similarly uses Merkle roots to commit to chunk data within shard blocks.
Merkle inclusion proofs are $O(\log n)$ and need only hashing to verify.
They give light clients no way to check encoding correctness or to sample: a proof shows that a piece belongs to the committed set but says nothing about whether the erasure coding was done correctly.
In both Polkadot and NEAR, encoding correctness is enforced through the validator set (validators who reconstruct and re-encode can detect errors), but this verification is not accessible to light clients or external observers.

% ----------------------------------------------------------
% III-C. VERIFICATION PARADIGMS
% ----------------------------------------------------------
\subsection{Verification Paradigms: DAS, VID, and Validator Recovery}
\label{subsec:verification}

The protocols differ most in how availability is verified.
The three paradigms below differ in their trust assumptions and in what a light client can check.

% --- DAS ---
\subsubsection{Data Availability Sampling (DAS)}

DAS, as formalized by Al-Bassam et al.~\cite{AlBassam_FraudDA}, enables light clients to verify data availability without downloading the full block.
A light client randomly requests $s$ coded pieces from the network; if all $s$ are returned, the client gains probabilistic confidence that the full data is available.
Against a $(1{-}q)$-withholding adversary (one that has made fraction $q$ of the data unavailable), a client sampling $s$ independent pieces achieves confidence $1 - q^s$.

Celestia was the first production network to deploy DAS.
Its 2D RS structure means that light clients sample from the $2k \times 2k$ extended square.
Because an adversary must withhold at least ${\sim}25\%$ of the extended square to prevent reconstruction, each sample detects withholding with probability ${\ge}1/4$, so 16 random samples give a Celestia light client approximately $1 - (3/4)^{16} \approx 99\%$ confidence that the data is available, assuming that at least one honest full node is non-eclipsed and participates in the peer-to-peer network~\cite{AlBassam_LazyLedger}.

Avail implements DAS over its 1D RS-encoded rows with KZG commitments.
Avail light clients sample random cells from the data matrix; each cell is accompanied by a KZG opening proof.
With 10 random samples, an Avail light client achieves $1 - (1/2)^{10} \approx 99.9\%$ confidence ($1 - 1/1024$)~\cite{Avail_The_DA_Blockchain}.
Because each cell is checked against the committed row polynomial, the client also verifies encoding correctness without fraud proofs.
Avail light clients fetch cells from a Kademlia DHT run by light clients and fall back to full-node RPC; cells fetched over RPC are written back into the DHT, so later requests can be served by other light clients.

Ethereum activated DAS through PeerDAS (EIP-7594) in the Fusaka upgrade on 3~December 2025.
Under PeerDAS, each blob is extended from 64 to 128 evaluation points (columns), and each validator is responsible for custody of a subset of columns determined by its node ID.
In the Fulu specification every node custodies at least 4 of 128 custody groups (\texttt{CUSTODY\_REQUIREMENT}), nodes with attached validators custody at least 8 plus one per 32~ETH of balance (\texttt{VALIDATOR\_CUSTODY\_REQUIREMENT}), and every node samples 8 columns per slot (\texttt{SAMPLES\_PER\_SLOT})~\cite{EIP-7594}.
Each node samples only a few columns, so confidence comes from the aggregate: with custody spread over ${\sim}1{,}000{,}000$ validators, each of the 128 columns is very likely held by at least one honest validator.
PeerDAS also introduces a ``supernode'' role: nodes that subscribe to all 128~data column subnets and can reconstruct the full data matrix from partial columns. The Fulu networking specification requires at least one supernode on the network for reconstruction of missing data, an assumption that did not exist before PeerDAS, when every consensus node downloaded all blob data~\cite{EIP-7594}.

DAS on all three protocols relies on two assumptions: (i)~at least one honest, non-eclipsed full node must hold the complete data and respond to sampling requests, and (ii)~sufficient sampling participation must exist in the network to ensure that withholding is detected before the data propagation window closes.
An adversary capable of eclipsing a light client's peer connections can cause the client to falsely conclude that data is available (or unavailable) regardless of the true state.

% --- Validator-Set Recovery ---
\subsubsection{Validator-Set Recovery}

In Polkadot and NEAR the validator set stores and serves the erasure-coded data, and availability is attested through consensus instead of sampling.

Polkadot's ELVES protocol ties availability to finality~\cite{Jeff_Burdges_ELVES}.
After a parachain collator produces a block candidate, the relay chain validators erasure-code the Proof of Validity (PoV) and distribute one chunk per validator.
Each validator that successfully receives and stores its chunk signs an \emph{availability bitfield} (a bitmap indicating which PoV chunks it holds).
A PoV is considered available when $\geq 2/3$ of validators have signed its bitfield.
The GRANDPA finality gadget will not finalize a relay chain block unless the availability bitfields for all included parachain candidates have reached the $2/3$ threshold~\cite{Overview_of_polkadot, polkadot_article}.
Under the $2/3$ honest-majority assumption the guarantee is deterministic: unavailable data cannot be finalized.

NEAR's Nightshade sharding distributes chunk data across shard validators (chunk producers).
Each chunk producer is responsible for producing and serving the chunk parts for its assigned shard.
NEAR's availability enforcement relies on a ``lazy signer'' mechanism with probabilistic accountability: validators are required to attest to chunk availability, but enforcement is primarily through the consensus protocol's liveness requirements rather than explicit sampling~\cite{NEAR_Nightshade_Scaling_Blockchain}.
A challenge mechanism exists in the block structure for disputing invalid or unavailable chunks, but as of our observation period, this mechanism is ``not used today'' in practice.

Because availability evidence is a precondition for finality (Polkadot) or block production (NEAR), validators have a direct incentive to store and serve data.
Light clients, however, cannot check availability themselves: a Polkadot or NEAR light client must trust that the $\geq 2/3$ supermajority attested honestly, with only the validators' stake behind that trust.

% --- VID ---
\subsubsection{Verifiable Information Dispersal (VID)}

Espresso's Tiramisu replaces DAS with algebraic dispersal proofs~\cite{Espresso:Hotshot_and_Triamisu}.
It has three layers:

\squishlist
\item \textbf{Savoiardi (VID layer):} Data is encoded as a polynomial and evaluated at $n$ points (one per committee member). A vector commitment (using KZG) binds the evaluations to the polynomial commitment, producing a compact proof that each committee member received a correct evaluation. This provides \emph{bribery resistance}: even if an adversary bribes a threshold of committee members, the algebraic binding prevents them from claiming to hold data they do not possess.

\item \textbf{Mascarpone (certification layer):} Committee members who receive and verify their evaluations produce threshold signatures attesting to availability. A DA certificate is formed when $\geq 2/3$ of the committee has signed, providing a compact ($O(1)$-size) availability attestation.

\item \textbf{Cocoa (CDN layer):} A secondary retrieval path through an untrusted CDN that provides availability even if committee members go offline, with cryptographic proofs linking the fallback data to the original VID commitments.
\squishend

Confidence is deterministic: there is no sampling, and a block is finalized only with a DA certificate.
DAS detects withholding probabilistically, whereas Savoiardi's algebraic binding makes it infeasible for a committee member to claim falsely that it holds its piece.

Retrieval, however, requires cooperation from the committee or the Cocoa CDN.
There is no mechanism for an independent light client to sample data and verify availability without relying on the committee's attestations.
Espresso argues that DAS adds ``unnecessary redundancies'' given VID's algebraic guarantees~\cite{Espresso:Hotshot_and_Triamisu}; as a result, retrievability depends on the committee and the CDN.

% ----------------------------------------------------------
% CROSS-CUTTING: WITHHOLDING ENFORCEMENT
% ----------------------------------------------------------
\subsubsection{No Slashing for Data Withholding}

None of the six protocols slashes for pure data withholding.
Each has some way to detect or discourage it, but enforcement is always indirect:

\squishlist
\item \textbf{Polkadot:} Validators are slashed for casting invalid \emph{validity} votes (attesting that an invalid parachain block is valid), but there is no slashing for failing to store or serve erasure-coded chunks. A validator that withholds its chunk simply does not sign the availability bitfield, potentially delaying but not corrupting finality~\cite{Jeff_Burdges_ELVES}.

\item \textbf{Celestia:} Bad Encoding Fraud Proofs (BEFPs) allow any full node to prove that the block producer performed incorrect erasure coding~\cite{AlBassam_FraudDA}. However, BEFPs address encoding correctness, not data withholding. A block producer that correctly encodes data but subsequently withholds it faces no direct slashing penalty.

\item \textbf{Ethereum:} PeerDAS enforces availability through peer scoring and disconnection: validators that fail to serve their custody columns receive negative gossip scores and may be disconnected from the p2p network. There is no protocol-level slashing for custody failures~\cite{EIP-7594}.

\item \textbf{Espresso:} In the current Mainnet~0 deployment, enforcement is reputational only. No slashing mechanism is implemented for committee members that fail to serve their VID evaluations~\cite{Espresso:Hotshot_and_Triamisu}.

\item \textbf{NEAR:} A challenge mechanism exists in the block structure that would allow validators to dispute chunk unavailability, but this mechanism is ``not used today'' in the production network. Enforcement relies on the consensus protocol's liveness requirements~\cite{NEAR_Nightshade_Scaling_Blockchain}.

\item \textbf{Avail:} Slashing is implemented for equivocation (double-signing) but not for data withholding. Validators that fail to make sampled cells available face no direct economic penalty~\cite{Avail_The_DA_Blockchain}.
\squishend

Proving that a node \emph{has} data but \emph{refuses} to serve it is harder than proving that it served \emph{incorrect} data.
Withholding is an omission fault, and attributing an omission to malice instead of a network partition, latency or node failure is an open problem in distributed systems.

% ----------------------------------------------------------
% COMPREHENSIVE FEATURE COMPARISON TABLE
% ----------------------------------------------------------

\begin{table*}[t]
\centering
\caption{Feature and parameter comparison of the six DA protocols during the observation window.}
\label{tab:feature_comparison}
\setlength{\tabcolsep}{1.8pt}
\renewcommand{\arraystretch}{1.15}
\footnotesize
\begin{tabular}{l c c c c c c}
\toprule
\textbf{Feature} & \textbf{Polkadot} & \textbf{Ethereum} & \textbf{Celestia} & \textbf{Espresso} & \textbf{NEAR} & \textbf{Avail} \\
\midrule
Encoding &
  1D RS (validators) &
  1D RS (PeerDAS) &
  2D RS ($k{\to}2k$) &
  Poly.\ VID &
  1D RS (chunks) &
  1D RS (rows) \\
Commitment &
  Merkle trie &
  KZG &
  NMT &
  KZG + vector &
  Merkle &
  KZG \\
Redundancy &
  ${\sim}3\times$ &
  ${\sim}2\times$ &
  ${\sim}4\times$ &
  Configurable &
  ${\sim}3\times$ &
  ${\sim}2\times$ \\
DAS support &
  No &
  Yes (PeerDAS) &
  Yes (2D) &
  No (VID) &
  No &
  Yes (1D) \\
Light client DAS &
  No &
  Custody &
  Yes (16 samples) &
  No &
  No &
  Yes (10 samples) \\
Block time &
  6\,s (relay) &
  12\,s &
  6\,s &
  ${\sim}2$\,s &
  ${\sim}0.6$\,s &
  20\,s \\
Max DA capacity &
  10\,MiB$\times$100 &
  2.63\,MiB (BPO2) &
  8$\to$32\,MiB$^\ast$ &
  ${\sim}1$\,MiB &
  4\,MiB$\times$9 &
  4\,MiB \\
Protocol max MiB/s & $\sim$167 & $\sim$0.22 & 1.33$\to$5.31$^\ast$ & $\sim$0.37 & $\sim$59 & $\sim$0.21 \\
Finality &
  12--60\,s &
  ${\sim}13$\,min &
  ${\sim}6$\,s &
  ${\sim}2$\,s &
  ${\sim}2$\,s &
  20--60\,s \\
DA retention &
  25\,h &
  ${\sim}18$\,d &
  7\,d (sampling) &
  Node-set &
  ${\sim}1.5$\,d &
  Node-set$^\dagger$ \\
Fee model &
  Weight-based &
  Blob market &
  Gas-based &
  Per-byte &
  Gas-based &
  Weight-based \\
Withholding &
  Finality gate &
  Peer scoring &
  BEFPs (encoding) &
  Reputational &
  Unused challenge &
  None \\
Byzantine threshold &
  $f < n/3$ &
  $f < n/3$ &
  $f < n/3$ &
  $f < n/3$ &
  $f<n/3$/shard &
  $f < n/3$ \\
\bottomrule
\end{tabular}

\vspace{2pt}
\raggedright\footnotesize{$^\ast$Celestia's maximum square rose from 128 to 256 (8 to 32~MiB) on 3~Dec~2025. $^\dagger$Avail prunes state after 256 blocks, but block data is kept by default (Substrate \texttt{archive-canonical} block pruning).}
\end{table*}

Table~\ref{tab:feature_comparison} collects these parameters.
DAS protocols (Celestia, Avail, and Ethereum with PeerDAS) let light clients verify availability themselves but depend on network-level assumptions; validator-recovery (Polkadot, NEAR) and VID (Espresso) designs give deterministic guarantees that hold only if the validators or committee are honest.
The commitment scheme follows from the paradigm: KZG gives Ethereum, Avail and Espresso algebraic encoding checks, Celestia pairs sampling with NMTs and fraud proofs, and plain Merkle trees suffice where validators check the encoding themselves.
All six use the classical $f < n/3$ Byzantine threshold but apply it differently: Polkadot gates finality on it, Ethereum spreads custody, Celestia relies on sampling participation, Espresso on committee certification, NEAR on per-shard honesty, and Avail on sampling backed by algebraic commitments.
